\section{Introduction}\label{sec:intro}

Let $\FF$ denote $\R$ or $\C$.  A \emph{weighted design} of size $m$ and
rank $k$ consists of vectors $g_1,\dots,g_m\in\FF^k$ and weights
$t_1,\dots,t_m>0$ with $\sum_c t_c=1$ subject to the Parseval condition
\begin{equation}\label{eq:parseval}
  \sum_{c=1}^{m} t_c\, g_c^{\vphantom{*}} g_c^{*} \;=\; I_k .
\end{equation}
The condition says that the vectors resolve the identity on average.  The
selection problem asks whether $k$ of them resolve it outright: whether
there exists $C\subseteq\{1,\dots,m\}$ with $\lvert C\rvert=k$ and
\begin{equation}\label{eq:dominate}
  \SC \;:=\; \sum_{c\in C} g_c^{\vphantom{*}} g_c^{*} \;\psd\; I_k .
\end{equation}
Such a $C$ is said to \emph{dominate}.  We write $\Gtz(m,k)$ for the
assertion that every design of size $m$ and rank $k$ over $\R$ admits a
dominating $k$-subset, and $\GtzC(m,k)$ for its complex analogue.

At uniform weights the problem is classical.  Let $U\in\FF^{m\times k}$
satisfy $U^{*}U=I_k$, let $u_c^{*}$ be its $c$-th row, and put
$g_c:=\sqrt m\,u_c$, $t_c:=1/m$; then \eqref{eq:parseval} holds, and for
$U_C$ the $k\times k$ submatrix on the rows $C$ the requirement
\eqref{eq:dominate} reads
\[
  U_C^{*}U_C^{\vphantom{*}}\;\psd\;\tfrac1m I_k,
  \qquad\text{that is}\qquad
  \sigma_{\min}(U_C)\;\ge\;m^{-1/2},
  \qquad\text{that is}\qquad
  \lVert U_C^{-1}\rVert_2\;\le\;\sqrt m .
\]
That some $k\times k$ submatrix of a matrix with orthonormal columns has
inverse of spectral norm at most $\sqrt m$ is the conjecture of Goreinov,
Tyrtyshnikov and Zamarashkin \cite{GTZ1997}, raised there in the course of
bounding the error of pseudoskeleton approximation.  It is the exact
analytic content of the maximal-volume principle, and it remains open.

What is provable by that principle is the following \cite{GZT1997,GT2001}.
Among the finitely many $k$-subsets choose one of maximal volume, that is
one maximizing $\lvert\det U_C\rvert$.  It is nonsingular: by the
Cauchy--Binet formula
$\sum_{\lvert C\rvert=k}\lvert\det U_C\rvert^2=\det(U^{*}U)=1$, so the
volumes do not all vanish.  Put $B:=UU_C^{-1}$.  Replacing the
$j$-th row of $U_C$ by the $r$-th row of $U$ replaces that row by its
expansion $\sum_i B_{ri}(U_C)_i$ in the selected rows, and multilinearity
leaves one surviving term, so the new determinant is $B_{rj}\det U_C$.  An
entry of $B$ of modulus greater than $1$ would therefore name a $k$-subset
of strictly larger volume, and maximality gives $\lvert B_{rj}\rvert\le1$
throughout.  The rows of $B$ indexed by $C$ form $I_k$, whence
\[
  \tr\bigl((U_C^{*}U_C^{\vphantom{*}})^{-1}\bigr)
  \;=\;\lVert U_C^{-1}\rVert_F^2
  \;=\;\lVert B\rVert_F^2
  \;\le\;k+k(m-k),
\]
using $\lVert UX\rVert_F=\lVert X\rVert_F$ in the middle step.  Moreover
$U_C^{*}U_C^{\vphantom{*}}\preceq U^{*}U=I_k$, so
$(U_C^{*}U_C^{\vphantom{*}})^{-1}\psd I_k$; a Hermitian $k\times k$ matrix
dominating the identity has all its eigenvalues at least $1$ and so has
largest eigenvalue at most its trace diminished by $k-1$.  Therefore
$\lVert U_C^{-1}\rVert_2^2\le k(m-k)+1$ and
$\SC\psd\frac{m}{k(m-k)+1}I_k$.  The bound and the conjecture are related
by the identity
\begin{equation}\label{eq:deficit}
  k(m-k)+1\;=\;m\;+\;(k-1)(m-k-1),
\end{equation}
a ring identity in the substitution $k=a+1$, $m-k=b+1$.  The deficit
$(k-1)(m-k-1)$ vanishes precisely when $k=1$ or $m=k+1$, and is positive
otherwise; at $(m,k)=(7,3)$ it equals $6$, so that the maximal-volume
principle delivers $\sigma_{\min}^2\ge1/13$ where $1/7$ is asked.  Thus
the range in which the classical argument already suffices is exactly
\begin{equation}\label{eq:elementary}
  k\le1\quad\text{or}\quad m\le k+1 ,
\end{equation}
and every remaining case of the conjecture demands something the volume
count does not supply.

Beyond \eqref{eq:elementary} the known cases are few.  Nesterenko
\cite{Nesterenko2023} settled $(m,k)=(4,2)$ over $\R$ by an argument in
Pl\"ucker coordinates, and reformulated the rank-two problem as an
isoperimetric question for closed polygons \cite{Nesterenko2024}; Sengupta
and Pautov \cite{SenguptaPautov2026} then settled rank two over $\R$ for
every $m$.  The equality locus at rank two over $\R$ was described by
Nesterenko \cite{Nesterenko2026real}, and configurations attaining the
bound $1$ at every pair $(m,k)$ were constructed by him from the star
spaces of series-parallel graphs \cite{Nesterenko2025}.  Over $\C$ the
conjecture is false at rank two: at uniform weights some pair always
satisfies $\SC\psd\alpha I_2$ with $\alpha=2-2/\sqrt3<1$, and $\alpha$
cannot be raised, being attained when $4\mid m$ and approached otherwise
\cite{Nesterenko2026complex}.  The restricted-invertibility theorems
\cite{BourgainTzafriri1987,SpielmanSrivastava2012,MSS2022} select subsets
of size strictly below the rank with a least eigenvalue controlled by a
stable-rank ratio; the control degenerates as the size of the subset
approaches the rank, which is the whole of the present range.  Finally,
maximizing the volume is NP-hard and NP-hard to approximate
\cite{CivrilMagdonIsmail2009}, so the selection principle is an existence
statement and not an algorithm; on the algorithmic side one settles for
the polynomial factors of \cite{GuEisenstat1996}.

\medskip

This paper studies the problem with general weights.  Two things are
gained.  First, the weights become a coordinate: the vectors and the
weights may be varied independently, subject to \eqref{eq:parseval},
and the extremal configurations then appear not as isolated matrices but
as sections over the weight simplex.  Second, the substitution
$v_c:=\sqrt{t_c}\,g_c$ turns \eqref{eq:parseval} into the statement that
the matrix $V$ with rows $v_c^{*}$ is an isometry, and \eqref{eq:dominate}
into a comparison between a principal block of the projection $P=VV^{*}$
and the corresponding block of $\diag(t)$.  In those coordinates the
problem is a question about a compact set, which is what makes a
variational treatment possible.  For each fixed pair $(m,k)$ the weighted
statement is formally stronger than the classical one; over all sizes the
two are equivalent, since a design with rational weights is a uniform
design of larger size with repeated vectors, and arbitrary weights follow
by a limiting argument.

\medskip

We prove four things.

The first is a complete answer over $\C$.  For $k\ge2$ selection succeeds
in every complex design exactly when $m\le k+1$, and for $k\le1$ it
succeeds always (Theorem~\ref{thm:complex}).  Comparison with
\eqref{eq:elementary} is exact: over $\C$ the conjecture is true precisely
where the maximal-volume count already proves it, and false everywhere
else.  The complex problem is thereby closed, and with it the possibility
of a field-independent proof of the real conjecture.

The second is the accessible part of the real problem.  Rank two holds at
every size (Theorem~\ref{thm:ranktwo}), corank one and corank two hold at
every rank (Theorems~\ref{thm:corankone} and~\ref{thm:coranktwo}), and at
each fixed rank the problem is finite: if $\Gtz(m,k)$ holds for all
$m\le k(k+1)/2+1$ then it holds for all $m$ (Theorem~\ref{thm:window}).
The bound is one more than the dimension of the space of real symmetric
$k\times k$ matrices, and it is where a support-reduction argument applied
to \eqref{eq:parseval} begins to bite.  At rank three the surviving cells
are $(6,3)$ and $(7,3)$, and no symmetry of the problem reduces the pair
further: $(6,3)$ is self-dual and $(7,3)$ is carried by duality to
$(7,4)$, of corank three (Corollary~\ref{cor:twocells}).

The third is the structure of the equality locus.  The threshold in
\eqref{eq:dominate} is attained: at every size $m\ge k+1$ there are
designs in which some $k$-subset dominates and none dominates strictly
(Theorem~\ref{thm:tieseverywhere}).  At corank one the designs attaining
it are classified, and they form a section over the whole open weight
simplex, hence a set of dimension $m-1$ (Theorem~\ref{thm:corankone}).
The regular tetrahedron of Example~\ref{ex:tetra} is one of them, and at
$(5,3)$ there is another whose vectors are pairwise non-parallel
(Example~\ref{ex:diamond}), so the locus is not exhausted by the designs
obtained from smaller ones by splitting a vector.

The fourth is a variational principle for the real problem.  The space of
designs is not compact --- the weights range over an open simplex, and by
\eqref{eq:parseval} a vector diverges as its weight tends to zero --- but
in the coordinates $(P,t)$ it is the interior of a compact space on which
the objective $\max_{\lvert C\rvert=k}\lmin(W[C])$, $W=P-\diag(t)$,
extends continuously.  We show that the boundary of that compactification
is free: granting the same statement at $(m-1,k)$ and at $(m-1,k-1)$,
every point with a vanishing weight carries a dominating $k$-subset
(Theorem~\ref{thm:boundary}).  Under those two hypotheses, if $\Gtz(m,k)$
fails then the objective attains its minimum at an interior point, that
is at a genuine weighted design (Theorem~\ref{thm:variational}); at
$(6,3)$ both hypotheses are theorems of \S\ref{sec:real}, and at $(7,3)$
one of them is the cell $(6,3)$.  What remains of rank three is the
exclusion of critical points below the threshold at the two surviving
cells, and \S\ref{sec:interior} derives the system such a point satisfies:
a multiplier $\Lambda$ obeying a quadric law $g_c^{*}\Lambda g_c=v$ at
every $c$ together with a coverage law, and three exclusions cutting down
the possible configurations of active subsets.  The objective minimized in
\S\ref{sec:compactness} and the one differentiated in \S\ref{sec:interior}
are related by a congruence and need not share their minimizers; the
passage between them is isolated as a hypothesis
(Definition~\ref{def:extremalcex}) and carried by
Theorem~\ref{thm:rankthree} alone.

\medskip

One fact organizes everything.  Because the threshold is attained, and
attained on a set of full dimension in the weights, no proof of $\Gtz$ can
carry a uniform margin: for every $\varepsilon>0$ there are designs
admitting no $k$-subset with $\SC\psd(1+\varepsilon)I_k$
(Corollary~\ref{cor:nomargin}).  This is a constraint on the form of any
solution.  An induction
on the size or on the rank cannot transport a quantitative slack across a
step, because there is none to transport at the source.  A certificate for
\eqref{eq:dominate} --- a sum-of-squares identity, a dual multiplier, a
determinantal or interlacing bound --- must be exact on the extremal
locus, so that any strategy proving a strict inequality is refuted before
it is attempted.  A minimizing sequence cannot be handled term by term,
since the identity available at a degenerate limit fails at every positive
weight, however small (Theorem~\ref{thm:sharp}); this is why the
compactification, and not a density argument, is what the variational
principle needs.  The strategies excluded for this reason, and the exact
rational witnesses refuting them, are collected in
Appendix~\ref{app:graveyard}.

What is not proved is stated as such.  The two cells $(6,3)$ and $(7,3)$
remain open (Conjectures~\ref{conj:sixthree} and~\ref{conj:seventhree}),
and with them the whole of the real problem at rank three and above; the
sharp constant over $\C$ at rank three is bracketed but not determined
(Remark~\ref{rem:sharpc}).  The problems left open, and the shape a
solution must have, are collected in \S\ref{sec:open}.

\subsection*{Markers}

A theorem number carrying $\dagger$ is machine-verified;
\S\ref{sec:formalization} gives the correspondence.  A number carrying
$\ast$ rests on a conjecture, stated at that point and not proved here; a
hypothesis discharged inside the paper is carried in the statement and
does not trigger the marker.  Unmarked statements are proved in the
ordinary way.  Every numerical assertion below is exact rational or exact
algebraic arithmetic.

\subsection*{Notation}

For $x,y\in\FF^k$ we write $\langle x,y\rangle:=x^{*}y$, conjugate-linear
in the first argument and linear in the second, and
$\lVert x\rVert^2=\langle x,x\rangle$.  Throughout,
\[
  \ell_c:=\lVert g_c\rVert^2,\qquad s_c:=t_c\ell_c ;
\]
$\ell_c$ is the squared length of the $c$-th vector and $s_c$ its
\emph{leverage score}.  For a Hermitian matrix $A$ we write $A\psd0$ for
positive semidefiniteness and $A\pd0$ for positive definiteness, $A[C]$
for the principal submatrix on the index set $C$, $\lmin$ and $\lmax$ for
the extreme eigenvalues, and $\operatorname{spec}A$ for the spectrum.  A
\emph{$k$-subset} always means a subset of $\{1,\dots,m\}$ of cardinality
$k$.  The \emph{corank} of a cell $(m,k)$ is $m-k$.  A pair $(m,k)$ with
$m<k$ carries no design (Lemma~\ref{lem:span}), so every statement below
that quantifies over such a pair is empty.

\begin{figure}[ht]
  \centering
  \begin{tikzpicture}[
  x=5.35mm, y=5.35mm,
  every node/.style={font=\scriptsize, inner sep=1pt},
  holds/.style   ={fill=black!12, draw=black!45, line width=0.3pt},
  fails/.style   ={fill=black!58, draw=black!45, line width=0.3pt},
  reduces/.style ={pattern=north east lines, pattern color=black!60,
                   draw=black!45, line width=0.3pt},
  opencell/.style={fill=white, draw=black, line width=0.9pt},
  nodesign/.style={draw=black!25, line width=0.25pt,
                   dash pattern=on 0.5pt off 0.7pt}]

% the cell (m,k) is the unit square with lower left corner (m-1,k-1)
\newcommand{\bcell}[3]{\path[#3] (#1-1,#2-1) rectangle ++(1,1);}
\newcommand{\bfrontier}[2]{%
  \path[opencell] (#1-1,#2-1) rectangle ++(1,1);
  \fill[black] (#1-0.5,#2-0.5) circle[radius=0.75mm];}

% =======================  the real board  ============================
\begin{scope}
  \foreach \m in {1,...,9}{\bcell{\m}{1}{holds}}

  \bcell{1}{2}{nodesign}
  \foreach \m in {2,...,9}{\bcell{\m}{2}{holds}}

  \foreach \m in {1,2}{\bcell{\m}{3}{nodesign}}
  \foreach \m in {3,4,5}{\bcell{\m}{3}{holds}}
  \bfrontier{6}{3}
  \bfrontier{7}{3}
  \foreach \m in {8,9}{\bcell{\m}{3}{reduces}}

  \foreach \m in {1,2,3}{\bcell{\m}{4}{nodesign}}
  \foreach \m in {4,5,6}{\bcell{\m}{4}{holds}}
  \bcell{7}{4}{reduces}
  \foreach \m in {8,9}{\bcell{\m}{4}{opencell}}

  \foreach \m in {1,...,4}{\bcell{\m}{5}{nodesign}}
  \foreach \m in {5,6,7}{\bcell{\m}{5}{holds}}
  \bcell{8}{5}{reduces}
  \bcell{9}{5}{opencell}

  \foreach \m in {1,...,9}{\node[anchor=north] at (\m-0.5,-0.15) {\m};}
  \foreach \k in {1,...,5}{\node[anchor=east]  at (-0.15,\k-0.5) {\k};}
  \node[anchor=north] at (4.5,-1.00) {size $m$};
  \node[rotate=90, anchor=south] at (-1.00,2.5) {rank $k$};
  \node[anchor=south] at (4.5,5.20) {$\FF=\R$};
\end{scope}

% ======================  the complex board  ==========================
\begin{scope}[shift={(11.2,0)}]
  \foreach \m in {1,...,9}{\bcell{\m}{1}{holds}}

  \bcell{1}{2}{nodesign}
  \foreach \m in {2,3}{\bcell{\m}{2}{holds}}
  \foreach \m in {4,...,9}{\bcell{\m}{2}{fails}}

  \foreach \m in {1,2}{\bcell{\m}{3}{nodesign}}
  \foreach \m in {3,4}{\bcell{\m}{3}{holds}}
  \foreach \m in {5,...,9}{\bcell{\m}{3}{fails}}

  \foreach \m in {1,2,3}{\bcell{\m}{4}{nodesign}}
  \foreach \m in {4,5}{\bcell{\m}{4}{holds}}
  \foreach \m in {6,...,9}{\bcell{\m}{4}{fails}}

  \foreach \m in {1,...,4}{\bcell{\m}{5}{nodesign}}
  \foreach \m in {5,6}{\bcell{\m}{5}{holds}}
  \foreach \m in {7,8,9}{\bcell{\m}{5}{fails}}

  \foreach \m in {1,...,9}{\node[anchor=north] at (\m-0.5,-0.15) {\m};}
  \foreach \k in {1,...,5}{\node[anchor=east]  at (-0.15,\k-0.5) {\k};}
  \node[anchor=north] at (4.5,-1.00) {size $m$};
  \node[anchor=south] at (4.5,5.20) {$\FF=\C$};
\end{scope}

% ===========================  legend  ================================
\begin{scope}[shift={(0,-2.60)}]
  \path[holds]    (0,0) rectangle ++(0.85,0.6);
  \node[anchor=west] at (1.05,0.30) {proved};

  \path[opencell] (3.5,0) rectangle ++(0.85,0.6);
  \node[anchor=west] at (4.55,0.30) {open};

  \path[opencell] (6.7,0) rectangle ++(0.85,0.6);
  \fill[black] (7.125,0.30) circle[radius=0.75mm];
  \node[anchor=west] at (7.75,0.30) {open, rank three};
\end{scope}

\begin{scope}[shift={(0,-3.60)}]
  \path[reduces]  (0,0) rectangle ++(0.85,0.6);
  \node[anchor=west] at (1.05,0.30) {implied by the two open cells};

  \path[fails]    (8.6,0) rectangle ++(0.85,0.6);
  \node[anchor=west] at (9.65,0.30) {fails};

  \path[nodesign] (11.8,0) rectangle ++(0.85,0.6);
  \node[anchor=west] at (12.85,0.30) {no design};
\end{scope}

\end{tikzpicture}

  \caption{The decision board.  Rows are the rank $k$, columns the size
  $m$; cells with $m<k$ carry no design.  Over $\R$ the two open cells at
  rank three are $(6,3)$ and $(7,3)$; the hatched region reduces to them
  by Theorem~\ref{thm:window}.  Over $\C$ the answer is $m\le k+1$ for
  every $k\ge2$.}
  \label{fig:board}
\end{figure}
